\documentclass[../../../../SoftwareRequirementsSpecification.tex]{subfiles}
\begin{document}
% Richiede le macro definite in src/macros/useCaseMacros.tex
% (incluse nel preambolo di SoftwareRequirementsSpecification.tex)

\ucstart
\begin{xltabular}{\textwidth}{|l|X|}
  \hline
  \ucid{PT-02} \\ \hline
  \ucname{Registrazione nuovo Atleta} \\ \hline
  \ucactor{PT} \\ \hline
  \ucsecondaryactor{Atleta da Registrare} \\ \hline
  \ucdescription{Il PT crea l'account di un nuovo atleta e lo iscrive
    a una o più palestre in cui il PT lavora. Il sistema invita
    l'atleta ad accedere all'applicativo.} \\ \hline
  \ucprecond{Il PT deve essere autenticato. Il PT deve conoscere
    l'email dell'atleta da registrare. Il PT deve essere associato ad
    almeno una palestra (vedi Nota 2).} \\ \hline
  \ucpostcond{L'atleta ha un account nell'applicativo, iscritto ad
    almeno una palestra. Il sistema lo ha invitato ad accedervi. Il
    profilo dell'atleta resta incompleto finché l'atleta non esegue il
    caso d'uso A-01.} \\ \hline
  \ucmainflow{
    \item Il PT apre la pagina di creazione nuovo atleta.
    \item Il sistema mostra il form di inserimento dati dell'atleta
      (nome, cognome ed email) e l'elenco delle palestre a cui il PT è
      attualmente associato, con una casella di selezione per
      ciascuna.
    \item Il PT inserisce i dati e seleziona una o più palestre in cui
      l'atleta si allena.
    \item Il PT preme il pulsante "Crea atleta".
    \item Il sistema verifica che l'email sia valida, non sia già
      presente nel database e che sia selezionata almeno una
      palestra.
    \item Il sistema crea l'account dell'atleta, lo associa al PT,
      crea l'iscrizione dell'atleta a ciascuna palestra selezionata
      (vedi Nota 3) e genera una password temporanea.
    \item Il sistema mostra una pagina di conferma creazione.
    \item Il sistema invia un'email all'atleta con l'invito ad
      accedere all'applicativo, la password temporanea e le
      istruzioni per il primo accesso.
  } \\ \hline
  \ucaltflow{
    \textbf{5a. Email già presente} \newline
    - Il sistema mostra un messaggio di errore ("Atleta con questa
    email già presente"). \newline
    - Il flusso riprende dal passo 2. \newline
    \textbf{5b. Email non valida} \newline
    - Il sistema mostra un messaggio di errore ("Email non
    valida"). \newline
    - Il flusso riprende dal passo 2. \newline
    \textbf{5c. Nessuna palestra selezionata} \newline
    - Il sistema mostra un messaggio di errore ("Selezionare almeno
    una palestra"). \newline
    - Il flusso riprende dal passo 2.
  } \\ \hline
  \ucnote{
    \textbf{Nota 1:} il PT non sceglie la password dell'atleta: la
    genera il sistema. Il PT non la vede mai, perché l'email è
    inviata direttamente all'atleta. \newline
    La password temporanea serve solo al primo accesso. Con il caso
    d'uso U-01 l'atleta entra, e il sistema lo porta subito al caso
    d'uso A-01, dove sceglie una propria password e completa il
    profilo. Finché non l'ha fatto, l'account resta inutilizzabile
    per qualsiasi altra funzione. \newline
    \textbf{Nota 2:} se il PT non è associato a nessuna palestra, il
    sistema non mostra questo caso d'uso: il PT deve prima associarsi
    a una palestra con il caso d'uso PT-13. \newline
    \textbf{Nota 3:} l'iscrizione dell'atleta a una palestra (passo 6)
    è indipendente dal rapporto tra il PT e quella palestra: se in
    seguito il PT termina la propria associazione con la palestra
    (PT-15), l'iscrizione dell'atleta non viene toccata. Le iscrizioni
    a una palestra si creano solo in questo caso d'uso, al momento
    della registrazione dell'atleta: la specifica non prevede né un
    caso d'uso per aggiungere un'iscrizione a un atleta già
    registrato, né un modo per l'atleta di gestire da sé le proprie
    iscrizioni.
  } \\ \hline
\end{xltabular}
\end{document}
